  \documentclass{article}
\usepackage{amsmath}
\usepackage{graphicx}
\graphicspath{ {images/} }
\usepackage[spanish]{babel}
\usepackage{bbm}
\usepackage{amsthm}
\usepackage{authblk}
\usepackage{hyperref}
\usepackage[utf8]{inputenc}
\usepackage{mathrsfs}
\usepackage{amsfonts}
\usepackage{amssymb}
\usepackage{enumerate}
\usepackage[left=3cm,right=2cm,top=2cm,bottom=2cm]{geometry}
\usepackage{epsfig}
 \renewcommand{\baselinestretch}{0.93}
 \thispagestyle{empty}
 \pagestyle{empty}
\setlength{\textheight}{53pc} \setlength{\textwidth} {36pc}
\setlength{\topmargin}{-4mm}
\usepackage{multicol}
\newcommand*{\email}[1]{%
    \normalsize\href{mailto:#1}{#1}\par
    }
\setlength{\columnsep}{1cm}
\newcommand{\N}{\mathbb{N}}
\newcommand{\calM}{\cal{M}}
\newcommand{\calP}{\cal{P}}
\newcommand{\F}{\mathbb{F}}
\newcommand{\R}{\mathbb{R}}
\newcommand{\Z}{\mathbb{Z}}
\newcommand{\Q}{\mathbb{Q}}
\newcommand{\C}{\mathbb{C}}
\newcommand{\bbH}{\mathbb{H}}


\theoremstyle{definition}
\newtheorem{ejer}{Ejercicio}
\author{Marcos Morales}
\affil{Universidad Finis Terrae}
\title{Primera Semana\\}



	


\begin{document}


\begin{picture}(400, 75)
   \put(-40,15){\includegraphics[width=5cm]{UFT.png}}
\end{picture}


\begin{center}
\huge{Primera Semana}
\end{center}



\begin{center}
	\large{Cálculo Vectorial}
\end{center}



\begin{center}
\setlength{\parindent}{0cm}\large{Clases centradas para capacitar a los estudiantes de ingeniería en Cálculo Vectorial. \\}
\end{center}




\vspace{1cm}

\begin{center}
	\large{Contenidos:}
\end{center}

\setlength{\parindent}{2.9cm}- El espacio $\mathbb{R}^n$ \\
\phantom{abefwfewefwfefffw} - Operaciones entre vectores	\\
\phantom{abefwfewefwfefffw} - Producto punto y norma 	\\
\phantom{abefwfewefwfefffw} - Proyecciones y producto cruz\\
\phantom{abefwfewefwfefffw} - Rectas en el espacio\\
\phantom{abefwfewefwfefffw} - Ecuación del plano\\



\vspace{6cm}


\begin{center}
\today
\end{center}
\begin{center}
Santiago, Chile
\end{center}


\pagestyle{empty}


\setcounter{page}{1}
\pagestyle{plain}
\setlength{\parindent}{0cm}





\newpage



\begin{center}
\textbf{\large{El espacio $\mathbb{R}^n$}} \\
\end{center}

Para propósitos de este curso entenderemos el espacio $\mathbb{R}^n$ como el espacio de todas las tuplas de la forma $(x_1,x_2,..., x_n)$ tal que cada coordenada es un número real, en otras palabras

\begin{center}
$\mathbb{R}^n = \lbrace (x_1,x_2,..., x_n) : x_1,x_2,... ,x_n \in \mathbb{R} \rbrace$
\end{center}

A lo elementos de $\mathbb{R}^n$ les llamaremos \textbf{vectores} y los denotaremos con una letra y una flecha, por ejemplo $\overrightarrow{u} = (1,2,\sqrt{2})$ o bien $\overrightarrow{w} = (5,-4,75)$. \\

Se tiene que $\mathbb{R}^2$ se identifica como el plano cartesiano mientras que $ \mathbb{R}^3$ se identifica como el espacio. Además $\mathbb{R}^1$ se puede identificar como el mismo conjunto de los número reales, es decir $\mathbb{R}^1 =\mathbb{R}$. \\



Para los que hayan visto álgebra lineal se tiene que $\mathbb{R}^n$ es un espacio vectorial de dimensión $n$ y por lo tanto tiene sentido llamar vectores a sus elementos, porque por definición un vector es un elemento de un espacio vectorial. \\



\begin{center}
\textbf{\large{Operaciones entre vectores}} \\
\end{center}


Sean dos vecotres $\overrightarrow{u},\overrightarrow{v} \in \mathbb{R}^n$, donde 

\begin{center}
$\overrightarrow{u} = (u_1,u_2,... ,u_n)$
\end{center}

\begin{center}
$\overrightarrow{v} = (v_1,v_2,... ,v_n)$
\end{center}


Podemos definir la \textbf{suma de vectores } como


\begin{center}
$\overrightarrow{u}+ \overrightarrow{v} = (u_1+v_1,u_2+v_2,... ,u_n+v_n) $
\end{center}

Es decir, la suma coordenada a coordenada, tenemos que la suma cumple con todas las propiedades de la suma en los número reales   \\

- $\overrightarrow{u}+ \overrightarrow{v} = \overrightarrow{v}+ \overrightarrow{u}$  (Conmutatividad) \\

- $ (\overrightarrow{u}+ \overrightarrow{v}) + \overrightarrow{w}  = \overrightarrow{u}+ (\overrightarrow{v} +\overrightarrow{w})$ (Asociatividad) \\

- Existe un elemento neutro $\overrightarrow{0}=(0, 0, ... ,0)$, tal que $\overrightarrow{u}+ \overrightarrow{0} = \overrightarrow{0}+ \overrightarrow{u}= \overrightarrow{u}+ \overrightarrow{u}$ para todo $\overrightarrow{u} \in \mathbb{R}^n$ (Existencia del elemento neutro)\\

- Para todo $\overrightarrow{u} \in \mathbb{R}^n$ existe un elemento $\overrightarrow{v} \in \mathbb{R}^n$ tal que $\overrightarrow{u}+ \overrightarrow{v}= \overrightarrow{0}$	(Existencia del elemento inverso)\\

Al inverso aditivo se le denota por $-\overrightarrow{u}$. Esta suma es una extensión de la suma habitual en $\mathbb{R}$ al espacio $\mathbb{R}^n$. \\

\newpage


Sean $\alpha \in \mathbb{R}$ y $\overrightarrow{u} \in \mathbb{R}^n$, donde $\overrightarrow{u} =(u_1,u_2,... ,u_n)$ podemos definir la \textbf{multiplicación por escalar} como

\begin{center}
$\alpha\overrightarrow{u} = (\alpha u_1,\alpha u_2,... ,\alpha u_n)$
\end{center}

Dados $\alpha,\beta \in \mathbb{R}$ y $\overrightarrow{u},\overrightarrow{v} \in \mathbb{R}^n$. Se cumple que

\begin{center}
$\alpha(\beta \overrightarrow{u}) = (\alpha \beta)\overrightarrow{u} = \beta(\alpha \overrightarrow{u})$ (Asociatividad)
\end{center}

\begin{center}
$\alpha(\overrightarrow{u}+\overrightarrow{v})=\alpha\overrightarrow{u} + \alpha \overrightarrow{u}$ (Distributividad)
\end{center}

\begin{center}
$(\alpha+\beta)\overrightarrow{u} = \alpha \overrightarrow{u} + \beta \overrightarrow{u}$ (Distributividad)
\end{center}

\begin{center}
$ 0 \overrightarrow{u} = \alpha \overrightarrow{0} = \overrightarrow{0}$
\end{center}

\begin{center}
$ (-1)\overrightarrow{u} = -\overrightarrow{u} $ 
\end{center}


\textbf{Ejemplo 1: } Sean $\overrightarrow{u}=(1,-3,4)$, $\overrightarrow{v}=(-1,2,-12)$ y $\overrightarrow{w} = (2,4,-17)$. Entonces

\begin{align*}
2\overrightarrow{u}-5\overrightarrow{v}+7\overrightarrow{w} &= 2(1,-3,4)-5(-1,2,12)+7(2,4,-17) \\
&= (2,-6,8)+(5,-10,-60)+(14,28,-119) \\
&= (21,12,-171) \\
\end{align*}


\textbf{Ejemplo 2: } Sean $\overrightarrow{u}=(1, 3, 5)$, $\overrightarrow{v}=(-8,4,3)$. Entonces

\begin{align*}
-4(\overrightarrow{u}+\overrightarrow{v}) &= -4( (1,3,5)+(-8,4,3)) \\
&= -4(-7,7,8) \\
&= (28,-28,-32) \\
\end{align*}





\newpage


\begin{center}
\textbf{\large{Producto punto, norma y desigualdad de Cauchy-Schwartz}} \\
\end{center}

Sean dos vecotres $\overrightarrow{u},\overrightarrow{v} \in \mathbb{R}^n$, donde 

\begin{center}
$\overrightarrow{u} = (u_1,u_2,... ,u_n)$
\end{center}

\begin{center}
$\overrightarrow{v} = (v_1,v_2,... ,v_n)$
\end{center}


Podemos definir el \textbf{ producto punto } como



\begin{center}
$\overrightarrow{u} \cdot \overrightarrow{v} = u_1v_1+u_2v_2+... +u_nv_n$
\end{center}


Notemos que el producto punto da como resultado un número real, es decir el producto punto es un operador que toma dos vectores y los lleva a un número real, esto se puedo anotar como $\cdot : \mathbb{R}^n \times \mathbb{R}^n \to \mathbb{R}$. \\

Sea $\overrightarrow{u} \in \mathbb{R}$ se puede definir la \textbf{norma}  $||.|| :\mathbb{R}^n \to \mathbb{R}$ como

\begin{center}
$|| \overrightarrow{u} || = \sqrt{u_1^2+u_2^2 +... +u_n^2}$
\end{center}

En el caso de $\mathbb{R}^2$ y $\mathbb{R}^3$ la norma se entiende como la distancia del vector al origen (el vector $\overrightarrow{0}$) mientras que el producto punto se entiende como el prducto entre la norma de un vector por la proyección sobre otro vector. En general se tiene que si $\alpha$ es el ángulo entre los vectores $\overrightarrow{u}$ y $\overrightarrow{v}$, entonces

\begin{center}
$\overrightarrow{u} \cdot \overrightarrow{v} = ||\overrightarrow{u}|| ||\overrightarrow{v}|| \cos(\alpha)$
\end{center}

Esta última igualdad fue demostrada en clases a través del teorema del coseno. Notemos que esto también nos permiete extender el concepto de ángulo entre vectores a dimensiones más grandes, en general se tiene que el ángluo entre dos vectores viene dado por

\begin{align*}
\displaystyle \alpha = \arccos \left( \frac{\overrightarrow{u} \cdot \overrightarrow{v}}{||\overrightarrow{u}|| || \overrightarrow{v}||} \right) 
\end{align*}

Mientras que la distancia entre los dos vectores viene dada por

\begin{center}
$d(\overrightarrow{u} ,\overrightarrow{v} ) = || \overrightarrow{u} -\overrightarrow{v} || $
\end{center}

\textbf{Ejemplo 1: } Sean $\overrightarrow{u} = (5,-2,1)$ y $\overrightarrow{v} = (1,-7,-1) $, luego el producto punto entre los vectores viene dado por

\begin{center}
$\overrightarrow{u} \cdot \overrightarrow{v} = 5 \cdot 1 + -2 \cdot -7 + 1 \cdot -1 = 5+14-1 =18 $
\end{center}


\textbf{Ejemplo 2: } Sean $\overrightarrow{u} = (5,-2,1)$ y $\overrightarrow{v} = (1,-7,-1) $, luego la distancia entre los dos vectores viene dado por


\begin{align*}
d(\overrightarrow{u} ,\overrightarrow{v} ) &= || \overrightarrow{u} -\overrightarrow{v} || \\
& ||(5,-2,1)-(1,-7,-1)|| \\
&= || (4,5,2)|| \\
&= \sqrt{4^2+5^2+2^2}\\
&= \sqrt{45} \\
&= 3\sqrt{5} 
\end{align*}

\newpage



\textbf{Ejemplo 3: } Sean $\overrightarrow{u} = (4,6,3)$ y $\overrightarrow{v} = (1,1,-1) $, luego el ángulo entre estos dos vectores viene dado por


\begin{align*}
\displaystyle \alpha &= \arccos\left( \frac{\overrightarrow{u} \cdot \overrightarrow{v}}{|| \overrightarrow{u} || ||\overrightarrow{v}||} \right)  \\
&= \arccos\left( \frac{4+6-3}{\sqrt{61} \sqrt{3}} \right)  \\
&= \arccos\left( \frac{7}{\sqrt{183}} \right)  \\
\end{align*}


\textbf{La norma y el producto punto cumplen con las siguientes propiedades}


\begin{center}
$||\overrightarrow{u} ||^2 = \overrightarrow{u}  \cdot \overrightarrow{u}  $
\end{center}

\begin{center}
$||\alpha \overrightarrow{u} || = |\alpha| || \overrightarrow{u} ||$
\end{center}

\begin{center}
$|| \overrightarrow{u} || \geq 0$ 
\end{center}

\begin{center}
$|| \overrightarrow{u} || = 0 \Leftrightarrow \overrightarrow{u} =\overrightarrow{0} $ 
\end{center}

\begin{center}
$|| \overrightarrow{u} + \overrightarrow{v}|| \leq ||\overrightarrow{u}|| +||\overrightarrow{v}|| $ (Desigualdad Triangular) 
\end{center}


\begin{center}
$ \overrightarrow{u} \cdot ( \overrightarrow{v} + \overrightarrow{w}) = \overrightarrow{u} \cdot \overrightarrow{v} + \overrightarrow{u} \cdot \overrightarrow{w}$ (Distributividad) 
\end{center}

\begin{center}
$\overrightarrow{u} \cdot \overrightarrow{u} \geq 0$
\end{center}


\begin{center}
$\overrightarrow{u} \cdot \overrightarrow{v} = \overrightarrow{v} \cdot \overrightarrow{u}$ (Conmutatividad)
\end{center}

\begin{center}
$\alpha (\overrightarrow{u} \cdot \overrightarrow{v}) = (\alpha \overrightarrow{u}) \cdot \overrightarrow{v} = \overrightarrow{u} \cdot (\alpha \overrightarrow{v})$ (Asociatividad)
\end{center}






\begin{center}
$\overrightarrow{u} \cdot \overrightarrow{v} \leq ||\overrightarrow{u}|| ||\overrightarrow{v}||$ (Desigualdad de Cauchy-Schwartz)
\end{center}



\textbf{Definición: } Dados dos vectores $\overrightarrow{u}, \overrightarrow{v} \in \mathbb{R}^n$ se dice que son \textbf{ortogonales } o \textbf{perpendiculares} si y sólo sí $\overrightarrow{u} \cdot \overrightarrow{v} = 0$ \\

\textbf{Ejemplo: } $\overrightarrow{u}= (-1,1,1)$ y $\overrightarrow{v} = (2,1,1)$, entonces

\begin{center}
$\overrightarrow{u} \cdot \overrightarrow{v} = -2+1+1=0 $
\end{center}


Luego $\overrightarrow{u}$ y $\overrightarrow{u}$ son perpendiculares




\newpage


\begin{center}
\textbf{\large{Proyecciones y producto cruz}} \\
\end{center}


Dados dos vectores $\overrightarrow{u}$ y $\overrightarrow{v}$ se define la proyección de $\overrightarrow{u}$ sobre $\overrightarrow{v}$ como

\begin{center}
$\displaystyle \text{Proy}_{\overrightarrow{u} \to \overrightarrow{v}} = \left( \frac{\overrightarrow{u} \cdot \overrightarrow{v}}{|| \overrightarrow{v}||^2} \right)  \overrightarrow{v}$
\end{center}

Notemos que $\displaystyle \frac{\overrightarrow{u} \cdot \overrightarrow{v}}{|| \overrightarrow{v}||^2}$ es un número real. \\


\textbf{Ejemplo: } Tomando $\overrightarrow{u} =(2,5)$ y $\overrightarrow{v} =(3,4) $, entonces

\begin{align*}
\displaystyle \text{Proy}_{\overrightarrow{u} \to \overrightarrow{v}} &= \left( \frac{\overrightarrow{u} \cdot \overrightarrow{v}}{|| \overrightarrow{v}||^2} \right)  \overrightarrow{v} \\
&= \left( \frac{(2,5) \cdot (3,4)}{||(3,4)||^2} \right)  (3,4) \\
&= \frac{6+20}{9+16}   (3,4) \\
&= \frac{26}{25}   (3,4) \\
&= \left( \frac{78}{25},\frac{104}{25} \right) 
\end{align*}


En $\mathbb{R}^3$, podemos escribir $\hat{\imath}=(1,0,0)$, $\hat{\jmath}=(0,1,0)$ y $\hat{k}=(0,0,1)$, de esta forma todo vector $\overrightarrow{u} =(u_1,u_2,u_3)$ se puede escribir de la forma

\begin{center}
$\overrightarrow{u}= u_1 \hat{\imath} +u_2\hat{\jmath}+u_3\hat{k}$
\end{center}




\textbf{Definición: } Dados dos vectores $\overrightarrow{u}$ y $\overrightarrow{u}$ en $\mathbb{R}^3$, surge la necesidad de determinar un vector $\overrightarrow{w}$ que sea perpendicular tanto a $\overrightarrow{u}$ como a $\overrightarrow{w}$. Esta definición será escrita sólo para $\mathbb{R}^3$, definimos entonces el \textbf{producto cruz} como

\begin{center}
$ \overrightarrow{u} \times \overrightarrow{v} = \begin{vmatrix}
i & j & k\\
u_1 & u_2 & u_3\\
v_1 & v_2 & v_3
\end{vmatrix} =  \begin{vmatrix}
u_2 & u_3 \\
v_2 & v_3
\end{vmatrix} \hat{\imath}- \begin{vmatrix}
u_1 & u_3 \\
v_1 & v_3
\end{vmatrix} \hat{\jmath} + \begin{vmatrix}
u_1 & u_2 \\
v_1 & v_2
\end{vmatrix} \hat{k}$
\end{center}

Notemos que $\overrightarrow{u} \times \overrightarrow{v}$ es un vector que es perpendicular tanto a $\overrightarrow{u}$ como a $\overrightarrow{u}$, además se tiene que

\begin{center}
$||\overrightarrow{u} \times \overrightarrow{v} || =||\overrightarrow{u}||||\overrightarrow{v}|| \sen(\alpha)$
\end{center}


Donde $\alpha$ es el ángulo entre los dos vectores. \\


\textbf{Ejemplo: } Sean $\overrightarrow{u} =(1,1,1)$ y $\overrightarrow{v} = (1,-2,-1) $, entonces


\begin{center}
$ \overrightarrow{u} \times \overrightarrow{v} = \begin{vmatrix}
i & j & k\\
1 & 1 & 1\\
1 & -2 & -1
\end{vmatrix} =  \begin{vmatrix}
1 & 1 \\
-2 & -1
\end{vmatrix} \hat{\imath}- \begin{vmatrix}
1 & 1 \\
1 & -1
\end{vmatrix} \hat{\jmath} + \begin{vmatrix}
1 & 1 \\
1 & -2
\end{vmatrix} \hat{k} = (1,2,-3)$
\end{center}


Dos vectores $\overrightarrow{u} $ y $\overrightarrow{v} $ se dicen \textbf{paralelos } si $\overrightarrow{u} = \lambda \overrightarrow{v} $ donde $\lambda \in \mathbb{R}$, se tiene que

\begin{center}
$\overrightarrow{u} $ es paralelo a $\overrightarrow{v}$ si y sólo sí $\overrightarrow{u} \times \overrightarrow{v} = \overrightarrow{0}  $
\end{center} 

\newpage

El producto cruz cumple con las siguientes propiedades \\

- $\overrightarrow{u} \times \overrightarrow{v} = - \overrightarrow{v} \times \overrightarrow{u} $ \\


- $(\overrightarrow{u} \times \overrightarrow{v}) \times \overrightarrow{z} = \overrightarrow{u} \times ( \overrightarrow{v} \times \overrightarrow{z} ) $ (Asociatividad) \\

- $\alpha (\overrightarrow{u} \times \overrightarrow{v}) = (\alpha \overrightarrow{u}) \times \overrightarrow{v}$ \\

-$(\overrightarrow{u}+\overrightarrow{v}) \times \overrightarrow{z} = \overrightarrow{u} \times \overrightarrow{z} +\overrightarrow{v} \times \overrightarrow{z}$ \\





\begin{center}
\textbf{\large{Vectores Unitarios}} \\
\end{center}


\textbf{Definición:} un vector $u$ se dice unitario si $|| u || =1$ \\

por ejemplo el vector $(1,0)$ es unitario. Dado un vector cualquiera $v$, se puede crear un vector unitario simplemente dividiendo por la norma, generando el vector

\begin{align*}
    \frac{v}{|| v||}
\end{align*}

Este nuevo vector tiene la misma dirección que $v$, pero es un vector unitario, a este proceso se le llama \textbf{normalizar un vector}. Los vecotres unitarios $(1,0,0)$, $(0,1,0)$ y $(0,0,1)$, también llamdos vectores canónicos, son muy útiles y se denotan de la siguiente manera

\begin{align*}
    \hat{\imath} &= (1,0,0) \\
    \hat{\jmath} &= (0,1,0) \\
    \hat{k} &= (0,0,1)
\end{align*}

Cualquier vector $(a,b,c)$ se puede expresar como

\begin{align*}
    (a,b,c) = a\hat{\imath}+b\hat{\jmath} +c\hat{k}
\end{align*}










\newpage

\begin{center}
\textbf{\large{Rectas en el espacio}} \\
\end{center}


Una recta en el espacio se puede identificar a través de un vector director $\overrightarrow{d}$ y un vector posición $\overrightarrow{p}$, de esta forma una recta $L$ viene dada por

\begin{center}
$L = \lbrace \overrightarrow{p} + \lambda \overrightarrow{d} : \lambda \in \mathbb{R}\rbrace$
\end{center}


\begin{figure}[h]
    \centering
    \includegraphics[width=0.3\textwidth]{Imagenes/descarga.jpg}
    \caption{Visuaizaci\'on de la ecuaci\'on vectorial de una recta}
    \label{fig:etiqueta}
\end{figure}

Usualmente se denota como 


\begin{center}
$L : \overrightarrow{p} + \lambda \overrightarrow{d} $
\end{center}

Esta última es llamada \textbf{Ecuación vectorial de la recta}. Por ejemplo podemos considerar $\overrightarrow{d}=(1,2,5)$ y $\overrightarrow{p}=(-1,2,3)$ y entonces la recta determinada por estos vectores sería

\begin{center}
$L : (-1,2,3) + \lambda (1,2,5) $
\end{center}


Si hacemos variar $\lambda$ obtenemos todo el conjunto de puntos que pertenecen a esta recta, por ejemplo, con $\lambda=2$, obtenemos

\begin{center}
$(-1,2,3) + 2(1,2,5)= (-1,2,3) + (2,4,10)=(1,6,13)  $
\end{center}

Y por lo tanto el punto $(1,6,13)$ pertenece a la recta $L$. \\


Notemos que si consideramos $\overrightarrow{d}=(d_1,d_2,d_3)$ y $\overrightarrow{p}=(p_1,p_2,p_3)$, entonces

\begin{align*}
L &: \overrightarrow{p} + \lambda \overrightarrow{d} \\
&= (p_1,p_2,p_3)+\lambda(d_1,d_2,d_3) \\
&= (p_1+\lambda d_1,p_2 + \lambda d_2,p_3 +\lambda d_3)\\
\end{align*}

De esta forma obtenemos que las cooredenas $x$, $y$ y $z$ están dadas por


\begin{align*}
x&= p_1+\lambda d_1 \\
y&= p_2+\lambda d_2 \\
z&= p_3+\lambda d_3 \\
\end{align*}


Esta se conoce como \textbf{Ecuación paramétrica de la recta} y expresa a todas las coordenadas en función de un único parámetro $\lambda$. Notemos que con estas ecuaciones podemos detectar inmediatamente el vector director y el vector posición. En efecto, sea la recta $L$ dada por 


\begin{align*}
x&= 2+5\lambda  \\
y&= 3-8\lambda \\
z&= -4+\lambda  \\
\end{align*}


Podemos detectar que el vector director es $\overrightarrow{d}=(5,-8,1)$ y el vector posición es $\overrightarrow{p}=(2,3,-4)$. Notemos que si despejamos $\lambda$ en cada ecuación obtenemos que


\begin{align*}
\displaystyle \lambda = \frac{x-p_1}{d_1}  \\
\lambda = \frac{y-p_2}{d_2} \\
\lambda = \frac{z-p_3}{d_3} 
\end{align*}

De dónde se deduce que


\begin{center}
$\displaystyle \frac{x-p_1}{d_1} = \frac{y-p_2}{d_2}  = \frac{z-p_3}{d_3} $
\end{center}

Esto se conoce como \textbf{Ecuación general de la recta}. Podemos detectar inmediatamente el vector director y el vector posición si nos fijamos en los coeficientes que acompañan a las coordenadas. Por ejemplo, consideremos la recta dada por la ecuación general


\begin{center}
$\displaystyle \frac{x-5}{3} = \frac{y+2}{7}  = z+4 $
\end{center}

Entonces el vector director es $\overrightarrow{d}=(3,7,1)$ y el vector posición es $\overrightarrow{p}=(5,-2,-4)$. \\


\textbf{Distancia de un punto a una recta:} En clases se discutió y demostró que la distancia de un punto $A$ a una recta $L: \overrightarrow{p}+ \lambda \overrightarrow{d}$ viene dada por

\begin{center}
$\displaystyle D(A,L) = \frac{|| (A-\overrightarrow{p}) \times \overrightarrow{d}||}{|| \overrightarrow{d}||}$
\end{center}


\textbf{Ejemplo: } Calcular la distancia entre el punto $A=(1,3,5)$ y la recta $L$ dada por $\displaystyle \frac{x-1}{2} = \frac{y+2}{3}  = z $. \\

\textit{Respuesta: }Primero hay que encontrar el vector posición y el vector director de la recta, se tiene que $\overrightarrow{d}=(2,3,1)$ y  $\overrightarrow{p}=(1,-2,0)$. Entonces

\begin{align*}
\displaystyle D(A,L) &= \frac{|| ((1,3,5)-(1,-2,0)) \times (2,3,1)||}{|| (2,3,1) ||} \\
 &= \frac{|| (0,5,5) \times (2,3,1)||}{\sqrt{14}} \\
  &= \frac{|| (-10,10,10)||}{\sqrt{14}} \\
    &= \frac{10\sqrt{3}}{\sqrt{14}} \\
\end{align*}



\textbf{Propiedades entre dos rectas:} Si tenemos dos rectas $L_1: \overrightarrow{p}+ \lambda \overrightarrow{d}$  y  $L_2: \overrightarrow{u}+ \lambda \overrightarrow{v}$, con $\overrightarrow{d}=(d_1,d_2,d_3)$, $\overrightarrow{p}=(p_1,p_2,p_3)$, $\overrightarrow{u}=(u_1,u_2,u_3)$ y $\overrightarrow{v}=(v_1,v_2,v_3)$. Entonces hay dos posibilidades  \\

\textbf{Caso 1: }  $\displaystyle \frac{d_1}{v_1}= \frac{d_2}{v_2} = \frac{d_3}{v_3}$ (Los vectores directores son paralelos y pueden ser considerados iguales). Entonces hay dos posibilidades \\

- 1. $\displaystyle \frac{p_1-u_1}{d_1}= \frac{p_2-u_2}{d_2} = \frac{p_3-u_3}{d_3}$. En este caso las rectas son coincidentes (Es decir son la misma recta). \\

- 2. $\displaystyle \frac{p_1-u_1}{d_1}\neq \frac{p_2-u_2}{d_2} $ o bien $\displaystyle  \frac{p_2-u_2}{d_2} \neq \frac{p_3-u_3}{d_3}$. En este caso las rectas son paralelas \\ \\ \\





\textbf{Caso 2: }  $\displaystyle \frac{d_1}{v_1} \neq \frac{d_2}{v_2}$ o bien $\displaystyle \frac{d_2}{v_2} \neq \frac{d_3}{v_3}$ (Los vectores directores  no son paralelos). Entonces hay dos posibilidades \\


- 1. 

\begin{center}
$\displaystyle \begin{vmatrix}
p_1-u_1 & d_1 & v_1 \\
p_2-u_2 & d_2 & v_2 \\
p_3-u_3 & d_3 & v_3 \\
\end{vmatrix} = 0$
\end{center}

En este caso las rectas son secantes, es decir se intersectan y su ángulo de intersección corresponde al ángulo entre los vectores directores. \\



- 2. 

\begin{center}
$\displaystyle \begin{vmatrix}
p_1-u_1 & d_1 & v_1 \\
p_2-u_2 & d_2 & v_2 \\
p_3-u_3 & d_3 & v_3 \\
\end{vmatrix} \neq 0$
\end{center}

En este caso las rectas  no se intersectan y se dice que las rectas son sesgadas. \\


\textbf{Ejemplo:} Consideramos las rectas  $\displaystyle L_1 :  \frac{x-5}{3} = \frac{y+2}{7}  = z+4 $ y $\displaystyle L_2 :  \frac{x-5}{2} = \frac{y-2}{6}  = \frac{z-1}{4} $. Quremos ver que tipo de rectas tenemos. \\

\textit{Respuesta: } Tenemos que los vectores directores de cada recta son $\overrightarrow{d}= (3,7,1)$ y $\overrightarrow{v}=(2,6,4)$, como los vectores no son paralelos pues $\displaystyle \frac{3}{2} \neq \frac{7}{6}$. Entonces, es claro que estamos en el caso $2$. Los vectores posición vienen dados por $\overrightarrow{p}=(5,-2,-4)$ y $\overrightarrow{u}=(5,2,1)$. Notamos que

\begin{center}
$\displaystyle \begin{vmatrix}
0 & 3 & 2 \\
-4 & 7 & 6 \\
-5 & 1 & 4 \\
\end{vmatrix} = 20 \neq 0$
\end{center}


Por lo tanto las rectas son sesgadas

\newpage

\begin{center}
\textbf{\large{Ecuación del plano}} \\
\end{center}

Como se vió en clases un plano $\Pi$ esta determinado por dos vectores directores $\overrightarrow{d}$ y $\overrightarrow{v}$ y un vector posición $\overrightarrow{p}$ y se puede escribir de la forma

\begin{center}
$\displaystyle \Pi : \overrightarrow{p}+\lambda \overrightarrow{d} + t \overrightarrow{v}$ con $\lambda,t \in \mathbb{R}$
\end{center}


\begin{figure}[h]
    \centering
    \includegraphics[width=0.4\textwidth]{Imagenes/images.png}
    \caption{Visuaizaci\'on de la ecuaci\'on vectorial del plano}
    \label{fig:etiqueta}
\end{figure}



Es decir, al hacer variar $\lambda$ y $t$ obtenemos todos los puntos del plano $\Pi$, a esta ecuación se le conoce como \textbf{Ecuación vectorial del plano}. Entonces, se dice que un plano tiene dos grados de libertad mientras que una recta tiene sólo un grado de libertad. \\

Al vector $\overrightarrow{p}$ se le llama \textbf{vector posición} y a los vectores $\overrightarrow{d}$ y $\overrightarrow{v}$ se les llama \textbf{vectores directores}. \\ 

\textbf{Ejemplo: }Consideramos el plano $\Pi$ generado por los vectores $\overrightarrow{d}=(1,2,3)$ y $\overrightarrow{v}=(-3,4,7)$ y el vector posición $\overrightarrow{p}= (1,0,3)$. Entonces el plano $\Pi$ esta dado por

\begin{align*}
\displaystyle \Pi &: (1,0,3)+\lambda(1,2,3)+t(-3,4,7) \\
&: (1+\lambda-3t,2\lambda+4t, 3+3\lambda+7t)
\end{align*}


Al reemplazar por ejemplo $\lambda=1$ y $t=2$, obtenemos el punto $(-4,10,20)$, esto quier decir que $(-4,10,20) \in \Pi$. \\


Si bien la ecuación vectorial del plano describe un plano totalmente, no es muy cómoda, en muchas ocaciones es preferible otra descripción del plano. \\

Consideremos un vector $\overrightarrow{n}$ y un vector $\overrightarrow{p}$, podemos crear un plano dado por

\begin{center}
$\displaystyle \Pi = \lbrace (x,y,x)\in \mathbb{R}^3 : ((x,y,z)-\overrightarrow{p}) \cdot \overrightarrow{n}=0 \rbrace$
\end{center}


Esta constrcción determina un plano y además todo plano puede ser representado de esta forma. Si denotamos $\overrightarrow{p}=(p_1,p_2,p_3)$ y $\overrightarrow{n}=(n_1,n_2,n_3)$, entonces

\begin{center}
$((x,y,z)-\overrightarrow{p}) \cdot \overrightarrow{n}=0$
\end{center}

\begin{center}
$((x,y,z)-(p_1,p_2,p_3)) \cdot (n_1,n_2,n_3)=0$
\end{center}

\begin{center}
$(x-p_1,y-p_2,z-p_3) \cdot (n_1,n_2,n_3)=0$
\end{center}


\begin{center}
$n_1(x-p_1)+n_2(y-p_2)+n_3(z-p_3)=0$
\end{center}


Y concluimos que

\begin{center}
$n_1x+n_2y+n_3z-n_1p_1-n_2p_2-n_3p_3=0$
\end{center}

La cual se puede dejar de la forma 

\begin{center}
$Ax+By+Cz+D=0$
\end{center}

Donde $A=n_1$, $b=n_2$, $C=n_3$ y $D=-n_1p_1-n_2p_2-n_3p_3 $. A esta ecuación se le conoce como \textbf{Ecuación general del plano} y todo plano puede ser representado por esta ecuación. \\

Al vector  $\overrightarrow{n}=(A,B,C)$ se le llama \textbf{vector normal} al plano



\textbf{Ejemplo:} Calcular la ecuación del plano cuyo vector normal es $\overrightarrow{n}=(1,7,-3)$ y que pasa por el punto $(-3,2,4)$.\\

\textit{Respuesta:} Ya tenemos $A$, $B$ y $C$, estos vienen dados por las coordenadas del vector normal, es decir $A=1$, $B=7$ y $C=-3$, luego la ecuación general es de la forma

\begin{center}
$x+7y-3z+D=0$
\end{center}

Para determinar $D$, tenemos que reemplazar el punto $(-3,2,4)$ en la ecuación y despejar $D$, entonces


\begin{center}
$-3+7(2)-3(4)+D=0$
\end{center}

\begin{center}
$-3+14-12+D=0$
\end{center}

\begin{center}
$D=1$
\end{center}

Por lo tanto la ecuación del plano es


\begin{center}
$x+7y-3z+1=0$
\end{center}



\textbf{Propiedades entre dos planos:} Si tenemos dos planos $\Pi_1$  y  $\Pi_2$ con vectores normales $\overrightarrow{n}$ y $\overrightarrow{m}$ y vectores posición $\overrightarrow{p}$ y $\overrightarrow{u}$ respectivamente. Entonces hay dos posibilidades  \\

\textbf{Caso 1: }  $\displaystyle \frac{n_1}{m_1}= \frac{n_2}{m_2} = \frac{n_3}{m_3}$ (Los vectores normales son paralelos y pueden ser considerados iguales). Entonces hay dos posibilidades \\

- 1. $(\overrightarrow{p}-\overrightarrow{u}) \cdot \overrightarrow{n} =0$. En este caso los planos son coincidentes (son el mismo plano) \\

- 2. $(\overrightarrow{p}-\overrightarrow{u}) \cdot \overrightarrow{n} \neq 0$. En este caso los planos son paralelos (no se cortan) \\





\textbf{Caso 2: }  $\displaystyle \frac{n_1}{m_1} \neq \frac{n_2}{m_2}$ o bien $\displaystyle \frac{n_2}{m_2} \neq \frac{n_3}{m_3}$ (Los vectores normales  no son paralelos ). En este caso los planos se intersectan y el ángulo entre los dos planos viene dado por : \\

\begin{center}
$\displaystyle \alpha= \arccos\left( \frac{|\overrightarrow{n} \cdot \overrightarrow{m}|}{||\overrightarrow{n}|| || \overrightarrow{m}||}\right) $
\end{center}

A este ángulo se le conoce como \textbf{ángulo diedro} \\


\textbf{Distancia de un punto a un plano:} La distancia de un punto $A$ a un plano $\Pi$ con vector normal $\overrightarrow{n}$ y vector posición $\overrightarrow{p}$ está dada por

\begin{center}
$\displaystyle d(A,\Pi) = \frac{| (A-\overrightarrow{p}) \cdot \overrightarrow{n}|}{||\overrightarrow{n}||}$
\end{center}

Esto es claro porque es simplemente la proyección del vector $A-\overrightarrow{p}$ sobre el vector normal $\overrightarrow{n}$. \\

Con esta fórmula se pueden deducir todas las otras fórmulas, por ejemplo la distancias entre planos paralelos, distancia de una recta a un plano, etc. Pero hay que diferenciar todos los casos. \\





\end{document} 